\section{Selecció de texts de MuSeD}\label{sec:filtratge}

\subsection{Estudi pilot}

Abans de la selecció final, vam dur a terme un estudi pilot amb 6 voluntaris (3 homes i 3 dones), que van llegir 11 texts de MuSeD seleccionats perquè contenien humor o ironia.
El pilot va permetre identificar dues limitacions del disseny inicial: els texts eren massa llargs, ja que els participants tardaven uns 10 minuts a llegir-los tots, i calia refinar el protocol experimental (vegeu \autoref{sec:procediment}).
Aquestes observacions van motivar els criteris de selecció que descrivim a continuació, orientats a reduir la longitud dels texts i a garantir una durada raonable de la sessió.

\subsection{Criteris de selecció}

A partir del corpus MuSeD, es van seleccionar 40 texts (20 sexistes i 20 no sexistes) seguint els criteris següents:

\begin{enumerate}
    \item \textbf{Coherència entre modalitats}: només es van conservar les instàncies en què les etiquetes textual i multimodal coincidien (\texttt{sexist\_text} $=$ \texttt{sexist\_video}).
    \item \textbf{Acord total entre anotadors}: es van descartar les instàncies amb valors intermedis de \texttt{sexist\_soft} i només es van mantenir les que presentaven consens ($\texttt{sexist\_soft} \in \{0, 1\}$).
    \item \textbf{Longitud}: es van prioritzar els texts més curts o amb segments anotats més curts, per mantenir una durada raonable de la sessió experimental i facilitar l'anàlisi.
    Considerant una velocitat de lectura aproximada de 1.000 caràcters per minut, els 40 texts seleccionats suposen aproximadament 17 minuts de lectura. A aquest temps cal afegir-hi el calibratge i les entrevistes de \textit{stimulated recall} (vegeu \autoref{sec:procediment}).
    \item \textbf{Comprensibilitat autònoma}: es van descartar els texts amb vocabulari estrany per a un oient catalanoparlant.
    Quan era necessari, es van fer petites edicions per millorar-ne la llegibilitat, com ara afegir cometes o indicar qui parla en fragments dialogats.
    \item \textbf{Validació manual}: cada text candidat es va revisar individualment. Es va escoltar l'àudio original i es van inspeccionar visualment les anotacions de segments per confirmar-ne l'adequació (vegeu l'\autoref{sec:spans_viz}).
\end{enumerate}


Els texts sexistes tendeixen a ser més llargs que els no sexistes, i les anotacions de segments en cobreixen una proporció considerable.
Les taules~\ref{tab:corpus_stats} i~\ref{tab:corpus_labels} resumeixen les estadístiques descriptives i la distribució de categories dels texts seleccionats.
En total, els 40 texts contenen 2.901 paraules.
Observem que la llargada dels segments constitueix un alt percentatge de la llargada total dels texts.
La llista completa dels texts, amb les categories binàries assignades, es troba a l'\autoref{sec:appendix_texts}.

\iftables
\begin{table}[ht]
\small
    \centering
    \begin{tabular}{lcccc}
        \toprule & \multicolumn{2}{c}{No sexistes} &
          \multicolumn{2}{c}{Sexistes} \\
        \cmidrule(lr){2-3} \cmidrule(lr){4-5}
        & Mitjana & D.E. & Mitjana & D.E. \\
        \midrule
        Longitud (caràcters)      & 388,9 & 137,9 & 439,4 & 173,4 \\
        Nombre de segments  & 0,95  & 0,50  & 1,83  & 0,75  \\
        Longitud dels segments & 255,7 & 139,6 & 362,7 & 205,3 \\
        \bottomrule
    \end{tabular}
    \caption{Estadístiques descriptives dels 40 texts seleccionats (D.E.:
    desviació estàndard).}\label{tab:corpus_stats}
\end{table}
\fi

\iftables
\begin{table}[ht]
\small
    \centering
    \begin{tabular}{llc}
        \toprule & Categoria & $n$ \\
        \midrule
        \multirow{4}{*}{\rotatebox[origin=c,y=1cm]{90}{No sex.}}
        & Crítica/denúncia  & 11 \\
        & Ironia            &  1 \\
        & Humor             &  2 \\
        %& Sexisme reportat  &  0 \\
        \midrule
        \multirow{8}{*}{\rotatebox[origin=c]{90}{Sexistes}}
        & Estereotips       & 11 \\
        & Desigualtat       &  7 \\
        & Discriminació     &  6 \\
        & Sexisme implícit  &  2\textsuperscript{a} \\
        & Sexisme reportat  &  1 \\
        & Ironia            &  2\\
        & Humor             &  2\textsuperscript{b} \\
        & Objectificació    &  0 \\
        \bottomrule
    \end{tabular}
    \caption{Distribució de categories dels texts seleccionats.
    \textsuperscript{a}~Els 2 texts amb sexisme implícit també presenten estereotips.
    \textsuperscript{b}~Els 2 texts sexistes amb humor també presenten discriminació.
    }\label{tab:corpus_labels}
\end{table}
\fi

% TODO: triar figura
% \begin{figure}[ht]
%     \centering
%     \includegraphics[width=\linewidth]{figures/exemple_anotacio.pdf}
%     \caption{Exemple d'un text sexista amb els segments anotats ressaltats.
%     Cada fragment destacat correspon a una porció del text identificada com a
%     sexista pels anotadors del corpus MuSeD.}
%     \label{fig:exemple_anotacio}
% \end{figure}

\section{Procediment experimental}\label{sec:procediment}

% \section{Participants}
L'estudi va comptar amb 10 participants voluntaris (5 homes i 5 dones).
Tots eren estudiants o treballadors de la Universitat de Barcelona, majors d'edat i d'entre 22 i 28 anys.
Dos dels participants també havien estat anotadors del corpus MuSeD.

Abans de començar, es va informar cada participant dels objectius de l'estudi i es va obtenir el consentiment informat.
Se'ls va explicar que haurien de llegir 40 texts en castellà, extrets de vídeos de xarxes socials, i determinar si eren sexistes.
Primer es va calibrar el dispositiu perquè l'eye-tracker funcionés adequadament (vegeu l'\autoref{app:calibration}).
Els texts es van mostrar en ordre aleatori.
Els participants podien llegir cada text tantes vegades com fos necessari i, quan havien decidit la resposta, premien 0 per indicar ``no sexista'' i 1 per indicar ``sexista'' en un teclat d'ordinador.
Se'ls demanava que valoressin el contingut d'acord amb la seva pròpia visió del món i se'ls explicava que no hi havia una única resposta correcta, perquè la tasca era subjectiva.
Després, apareixia una pantalla en què havien d'indicar la confiança en la seva elecció: 1 (poca confiança), 2 (confiança mitjana) o 3 (molta confiança).
Abans de mostrar el text següent, apareixia una pantalla en blanc durant 500\,ms.
Els participants podien fer preguntes o corregir errors durant aquesta fase.

\subsection{\textit{Stimulated Recall}}

Un cop acabada la lectura, es feien entrevistes breus mentre es mostrava el procés de lectura d'alguns texts amb Tobii Pro Lab, on es podien veure les fixacions i les regressions.
La pantalla de l'ordinador i l'àudio de la conversa es van enregistrar amb OBS Studio.
Es van seleccionar 7 texts que representaven casos interessants.
Tot i que el corpus MuSeD disposa d'etiquetes de consens (\textit{gold standard}), les etiquetes individuals revelen que l'acord entre anotadors no sempre és absolut.
Per exemple, el text~4 té unanimitat en discriminació (1,0) però desacord en humor (0,67) i ironia (0,33); el text~30 mostra desacord en humor (0,67) i ironia (0,33); i el text~375 en desigualtat (0,67) i múltiples categories (0,33 cadascuna).
Els texts seleccionats van ser:
\begin{itemize}
    \item \textbf{Text 4} (sexista): ironia i humor amb desacord, discriminació unànime.
    \item \textbf{Text 30} (sexista): humor i ironia amb desacord, desigualtat.
    \item \textbf{Text 37} (no sexista segons \textit{gold standard}): estereotips i desigualtat, però ironia i humor amb desacord alt.
    \item \textbf{Text 107} (sexista): sexisme implícit amb desacord baix.
    \item \textbf{Text 166} (no sexista): crítica i humor sense desacord en categories principals.
    \item \textbf{Text 375} (sexista): múltiples categories amb desacord alt (ironia, sexisme implícit, estereotips, desigualtat).
    \item \textbf{Text 399} (no sexista): crítica, estereotips i sexisme reportat, tots amb desacord.
\end{itemize}

Es va utilitzar la llengua amb què cada participant se sentia més còmode per conversar sobre els texts: català o castellà.
Durant aquesta fase, se'ls preguntava si s'havien fixat en alguna paraula o frase per un motiu concret, o si havien pensat alguna cosa determinada sobre el text.
També se'ls demanava si havien trobat difícil alguna paraula o si els havia sorprès algun concepte.

\section{Models de llenguatge}

%\subsection{Models BERT adaptats}

En aquest treball utilitzem diversos models basats en l'arquitectura \textit{transformer}~\cite{vaswaniAttentionAllYou2017}, concretament models del tipus BERT (\textit{Bidirectional Encoder Representations from Transformers})~\cite{devlin2019bert}.
Són models \textit{encoder-only} que han assolit un bon rendiment en tasques de classificació de texts.
Usem mBERT, la versió multilingüe del BERT original; BETO~\cite{canete2023spanish}, un dels primers models adaptats específicament al castellà; i el més recent MrBERT-es~\cite{tamayo2026mrbert}, basat en ModerBERT.

El procés d'adaptació (\textit{fine-tuning}) consisteix a modificar aquests models preentrenats per a la tasca específica de detecció de sexisme.
La implementació inclou la preparació de les dades, la configuració del model i l'entrenament amb la biblioteca \textit{transformers} de Hugging Face\footnote{https://huggingface.co/docs/transformers/en/index}.
També vam entrenar els models només amb les dades que tenien acord total entre anotadors, però les diferències respecte dels models entrenats amb totes les dades van ser negligibles.

% \subsection{Models de llenguatge extensos}
%
% A més dels models adaptats, també experimentem amb models de llenguatge grans (LLMs) en configuracions de \textit{zero-shot} i \textit{few-shot}.
% Aquests enfocaments permeten avaluar la capacitat dels models per generalitzar sense necessitat d'un entrenament específic per a la nostra tasca.
%
% En \textit{zero-shot}, el model realitza la classificació sense cap exemple previ de la tasca.
% En \textit{few-shot}, es proporcionen uns quants exemples (normalment entre 1 i 5) per guiar la predicció del model.

\section{Mètodes d'interpretabilitat}\label{sec:metodes}

Per comprendre els patrons d'atenció dels models i comparar-los amb els patrons humans, fem servir la biblioteca Captum~\cite{captum}, que proporciona diversos mètodes d'atribució.
N'hem implementat quatre:

\begin{itemize}
    \item \textbf{Saliency}~\cite{simonyan2013deep}: mètode basat en gradients que calcula la derivada de la sortida respecte dels \textit{embeddings} d'entrada. És útil per identificar els tokens que tenen més impacte en la decisió del model.
    \item \textbf{Input $\times$ Gradient}~\cite{shrikumar2016not}: multiplica l'entrada pels gradients per obtenir una mesura de la importància de cada característica. Pot proporcionar atribucions més estables que el \textit{saliency} simple.
    \item \textbf{Integrated Gradients (IG)}~\cite{sundararajan2017axiomatic}: calcula la integral dels gradients al llarg d'un camí que va des d'una línia base, normalment el vector nul, fins a l'entrada real.
    També proporciona una mètrica de convergència (\textit{delta}) que indica la qualitat de l'atribució.
    \item \textbf{Layer-wise Relevance Propagation (LRP)}~\cite{LRP}: atribueix la predicció a les característiques d'entrada mitjançant la redistribució de puntuacions de rellevància des de la sortida fins a l'entrada.
    Per a una mateixa entrada, LRP produeix explicacions reproduïbles.
\end{itemize}

% \subsection{Implementació pràctica}
%
% El procés d'interpretabilitat segueix les següents etapes:
%
% \begin{enumerate}
%     \item Carregar el model entrenat i el tokenitzador corresponent.
%     \item Tokenitzar el text d'entrada i obtenir els \textit{embeddings}.
%     \item Aplicar el mètode d'atribució seleccionat.
%     \item Sumar les atribucions al llarg de la dimensió dels \textit{embeddings} per obtenir una puntuació per a cada token.
%     \item Visualitzar els resultats utilitzant codis de colors (verd per a contribucions positives, vermell per a negatives).
% \end{enumerate}
%
% El codi inclou funcions per generar visualitzacions HTML amb els tokens ressaltats segons la seva importància per a la predicció del model.

\section{Comparació amb patrons humans}

L'objectiu final d'aquest treball és comparar tres fonts d'informació: (i) els patrons obtinguts amb els mètodes d'interpretabilitat, (ii) els patrons humans registrats amb l'eye-tracker i (iii) les anotacions inicials del corpus.
Aquesta comparació permet avaluar si els models se centren en les mateixes parts dels texts que les persones, identificar possibles biaixos o limitacions dels models i explorar si les divergències entre l'atenció humana i la computacional es relacionen amb errors de classificació.

Cal tenir en compte que els tokens d'eye-tracking corresponen a paraules senceres, mentre que els models de llenguatge treballen amb tokens subparaula.
Quan una paraula es divideix en múltiples tokens, les atribucions d'aquests se sumen.

L'anàlisi comparativa es basa en les dades d'eye-tracking recollides amb el dispositiu Tobii, tal com es descriu a l'\autoref{sec:tobii}.
Els resultats s'analitzen quantitativament i qualitativament per caracteritzar la similitud entre els dos tipus d'atenció.

% \section{Disponibilitat del codi}

Tot el codi fet servir en aquest treball, incloent-hi la implementació dels models, els mètodes d'interpretabilitat, els scripts d'anàlisi i aquest treball escrit, estan disponibles públicament al repositori GitHub: \url{https://github.com/Pastells/eye\_tracker\_sexism}. 
